% !TeX root = ../main.tex
\chapter{System Requirement Study}

This chapter documents requirement analysis for Jan-Sunwai AI, including the baseline state of existing grievance systems, observed limitations, actor expectations, and formalized functional/non-functional requirements. The study was derived from implementation realities, municipal workflow assumptions, and validation through tests and user journeys [1, 2, 11, 12, 13].

\section{Existing System Overview}

Many conventional civic complaint channels depend on one or more of the following interfaces: physical counters, basic web forms, call center logging, or generic ticketing portals. While these methods can register complaints, they often suffer from weak input quality, manual routing dependency, and low traceability.

Key characteristics of existing systems observed in this problem context:

\begin{itemize}
\item Citizens are expected to manually choose departments, often without sufficient guidance.
\item Complaint narratives are text-heavy and unstructured, leading to inconsistent detail quality.
\item Photo evidence may be accepted but is rarely integrated into routing logic.
\item Status updates are frequently sparse or delayed.
\item Workflow metadata (who changed what and why) is often incomplete for audits.
\end{itemize}

\section{Limitations of the Existing System}

Seven recurring limitations were observed in legacy civic grievance systems: manual authority selection by citizens causing frequent misrouting; inconsistent complaint descriptions that require staff to rewrite or clarify cases before action; weak integration of visual evidence leaving photographs unused during triage; limited role-specific workflows that impose the same interface burden on very different personas; minimal lifecycle traceability making accountability reviews difficult; low transparency in queue behavior eroding citizen trust; and sparse administrative analytics that force planning decisions to rely on intuition rather than data [1, 2, 6].

\section{User Characteristics}

The system is designed for four operational personas with distinct interaction patterns and authorization boundaries.

The system serves four distinct operational personas. \textbf{Citizens} report issues and track outcomes: they upload images, review the AI-generated draft, submit their complaint, and monitor notifications and status updates. \textbf{Workers} resolve assigned field tasks by viewing their queue, managing availability, and marking completions with resolution notes. \textbf{Department Heads} manage departmental flow by filtering the department queue, updating complaint status, adding notes, and transferring or escalating cases when required. \textbf{Administrators} exercise end-to-end governance: approving workers, triaging ambiguous AI classifications, performing bulk operations, exporting reports, and monitoring analytics.

\section{Functional Requirements}

\begin{description}
\item[FR-01 (User Registration):] System shall allow self-registration for citizen and worker roles with validation and duplicate checks.
\item[FR-02 (User Authentication):] System shall establish secure cookie-session authentication on successful login and protect authorized routes, with JWT compatibility for API clients.
\item[FR-03 (Profile Management):] Authenticated users shall update selected profile fields such as full name and phone number.
\item[FR-04 (Password Recovery):] System shall support forgot-password and reset-password flows with time-bound tokens.
\item[FR-05 (Image Upload Validation):] Analyze endpoint shall accept valid JPG/PNG/WEBP images and reject unsupported or malformed payloads.
\item[FR-06 (AI-Assisted Analysis):] System shall analyze uploaded image and return category, confidence, location metadata, and draft status.
\item[FR-07 (Draft Generation):] System shall produce a complaint draft synchronously or via queued generation, with polling support.
\item[FR-08 (Complaint Creation):] System shall create a complaint record from reviewed analysis payload and user-provided description edits.
\item[FR-09 (Complaint Listing):] System shall provide role-filtered complaint lists with pagination and optional status filters.
\item[FR-10 (Complaint Detail View):] System shall provide complete complaint detail including status history and operational metadata.
\item[FR-11 (Status Update Workflow):] Department head/admin shall update complaint status with note and trigger citizen notification.
\item[FR-12 (Department Transfer):] Department head/admin shall transfer complaint ownership across departments when misrouted.
\item[FR-13 (Escalation):] Department head/admin shall escalate complaints to parent authorities when required.
\item[FR-14 (Worker Assignment):] System shall auto-assign eligible workers and support manual override by admin.
\item[FR-15 (Worker Operations):] Worker shall manage availability and mark assigned complaints as done.
\item[FR-16 (Notification Center):] System shall provide list, unread count, single-read, and mark-all-read operations.
\item[FR-17 (Triage Queue):] Admin shall review low-confidence complaints and approve/reject/correct labels.
\item[FR-18 (Analytics):] Admin/dept-head shall access overview and geospatial heatmap summaries.
\item[FR-19 (CSV Export):] Admin shall export complaint data to CSV for external reporting.
\item[FR-20 (Public Transparency Feed):] System shall expose anonymized complaint records without authentication.
\item[FR-21 (Health Monitoring):] System shall provide live, ready, model, and GPU health endpoints.
\item[FR-22 (API Version Alias):] System shall expose major endpoints under /api/v1 for integration stability.
\end{description}

\section{Non-Functional Requirements}

\begin{description}
\item[NFR-01 (Security):] Cookie-session auth (with JWT compatibility), RBAC, sanitization, secure headers, and account-safe reset flow shall be enforced.
\item[NFR-02 (Reliability):] API shall degrade gracefully under model or DB faults with explicit status semantics.
\item[NFR-03 (Performance):] Public and health endpoints shall respond consistently under mixed read-heavy load.
\item[NFR-04 (Usability):] Role-specific UI shall reduce cognitive load and support responsive layouts for desktop/mobile usage.
\item[NFR-05 (Maintainability):] Modular routers/services with documented runbooks shall support iterative enhancement.
\item[NFR-06 (Observability):] Process-time headers, rotating logs, and health routes shall support debugging and operations.
\item[NFR-07 (Scalability Readiness):] Queue-based generation and container deployment shall support controlled scaling.
\item[NFR-08 (Portability):] Project setup shall support Windows and Linux via scripts and Docker-based workflow.
\item[NFR-09 (Compliance Readiness):] Local-first inference and sanitized logs shall support stricter governance deployments.
\item[NFR-10 (Data Integrity):] Validation schemas and controlled update flows shall minimize malformed or inconsistent records.
\end{description}

\section{Hardware and Software Requirements}

\subsection{Hardware Requirements}

The baseline parameters established below define the minimum viable constraints required to support the proposed features.

\begin{itemize}
    \item \textbf{CPU:} Minimum 4-core modern processor; Recommended 6 to 8-core.
    \item \textbf{System RAM:} 12 GB minimum; 16 GB or higher recommended.
    \item \textbf{GPU VRAM:} 4 GB baseline (e.g., RTX 3050); 6 GB or higher recommended for faster inference.
    \item \textbf{Storage:} 10 GB free for models and environment; 20 GB for long-term audit logs and uploads.
    \item \textbf{Network:} Stable LAN/Internet for map tiling and package management.
\end{itemize}

\subsection{Software Requirements}

The baseline parameters established below define the minimum viable constraints required to support the proposed features [16, 11].

\begin{itemize}
    \item \textbf{Operating System:} Windows 10+ or modern Linux (Ubuntu 22.04+ recommended).
    \item \textbf{Python Runtime:} Version 3.10 or higher for backend logic and automation.
    \item \textbf{Node.js \& npm:} Current LTS range (v20+) for frontend build and utility tooling.
    \item \textbf{MongoDB:} Version 7.x class for document-oriented state persistence.
    \item \textbf{Ollama Runtime:} Local execution engine for vision and reasoning pipelines.
    \item \textbf{Docker Engine/Compose:} Version 24+ for containerized environment orchestration.
    \item \textbf{Web Browser:} Modern Chromium or Firefox class for role-specific dashboard access.
\end{itemize}

\section{Constraints}

Constraints in this project are not only technical limits but also governance and usability boundaries that directly influence system behavior. The platform must operate under security, role-separation, and infrastructure limitations common in civic environments, while still maintaining user clarity for citizens and workflow efficiency for officials. The following constraint groups define these implementation boundaries and informed several architecture and UI decisions.

\subsection{System Constraints and Security}

System constraints govern UI responsiveness, security boundaries, and operational limits to ensure robust public access. 
\textbf{UI and Interface:} Role-specific pages strictly enforce access boundaries. Interfaces must support mobile widths (375px baseline) and preserve data state during refresh cycles. Communication happens exclusively over authenticated REST endpoints via cookie sessions or JWT tokens, with strict \texttt{/api/v1} proxy routing.
\textbf{Hardware Limits:} Because local inference is VRAM-bound, AI response times fluctuate under load. Large file uploads are strictly limited by size and decode verifications to prevent resource exhaustion.
\textbf{Security and Governance:} To maintain citizen trust, status history and role checks are immutable [6, 23, 7]. The platform mandates strict input sanitization, MIME-type validations, robust token lifecycles (preventing account enumeration), and active CORS/security headers across all responses [7, 8, 9].

The system operates under several core assumptions and external dependencies. First, it assumes a reachable and healthy \textbf{MongoDB instance} for all read/write persistence. Second, the \textbf{Ollama runtime} must be correctly configured with the required vision and reasoning models to support analysis and draft generation. Third, accurate \textbf{environment variable configuration} is essential for stable authentication, CORS, and database mapping. Fourth, the system relies on \textbf{worker profile completeness} (department and service-area metadata) for the auto-assignment engine to function correctly. Finally, for production-style deployments, the host must support \textbf{Docker and Compose} to ensure environment reproducibility.

\section{Detailed Use-Case Narratives}

\begin{description}
\item[UC-01 (Citizen):] Register account, login, access dashboard $\rightarrow$ Duplicate username/email is rejected with validation message.
\item[UC-02 (Citizen):] Upload image and run analyze pipeline $\rightarrow$ Unsupported/corrupt image rejected with explicit reason.
\item[UC-03 (Citizen):] Edit draft and submit complaint $\rightarrow$ If analyze unavailable, user submits manual complaint text after retry guidance.
\item[UC-04 (Citizen):] Track status and receive notifications $\rightarrow$ If no notifications available, panel remains stable with unread count zero.
\item[UC-05 (Citizen):] Submit feedback after complaint resolution $\rightarrow$ Duplicate feedback attempt blocked with conflict response.
\item[UC-06 (Worker):] View assigned complaints and update availability $\rightarrow$ Worker pending approval cannot access worker-only operations.
\item[UC-07 (Worker):] Mark complaint done $\rightarrow$ If assignment mismatch occurs, operation is denied by role/ownership checks.
\item[UC-08 (Dept Head):] Filter department complaints and update status $\rightarrow$ Dept-head cannot update complaints outside own department.
\item[UC-09 (Dept Head):] Add notes and transfer/escalate complaint $\rightarrow$ Invalid complaint ID or missing target yields controlled validation error.
\item[UC-10 (Admin):] Approve/reject worker registrations $\rightarrow$ Already-approved/rejected states handled idempotently.
\item[UC-11 (Admin):] Operate triage queue and analytics $\rightarrow$ Empty triage queue still returns successful response.
\item[UC-12 (Public User):] View anonymized complaint board $\rightarrow$ No authenticated personal fields are exposed in response payload.
\end{description}

Requirements are classified using the MoSCoW methodology [20]. \textbf{Must-Have} features include user registration, authentication, complaint intake, status updates, auto-assignment, and health monitoring: these are essential for operational viability. \textbf{Should-Have} features cover authority escalation, worker status controls, the administrative triage queue, and analytics: these provide institutional-grade governance. \textbf{Could-Have} features such as CSV exports and the public transparency feed provide high-value reporting without blocking the core workflow. Finally, \textbf{Non-Functional criticals} including session security, sanitization, and observability are treated as mandatory baseline controls.

\section{Chapter Summary}

The requirement study establishes the transition from a manually intensive grievance flow to a structured, secure, and auditable digital workflow. Functional and non-functional requirements were formalized in implementation-aligned terms and mapped against realistic user roles and operational constraints.

